\documentclass[11pt]{article}
\usepackage[letterpaper,margin=0.85in]{geometry}
\usepackage[T1]{fontenc}
\usepackage{lmodern,microtype,amsmath,amssymb,booktabs,tabularx,array,longtable}
\usepackage{xcolor,enumitem,fancyhdr}
\usepackage[hyphens]{url}
\usepackage{hyperref}
\definecolor{navy}{HTML}{17324D}
\definecolor{teal}{HTML}{176B73}
\hypersetup{colorlinks=true,linkcolor=navy,citecolor=teal,urlcolor=teal,
 pdftitle={Binary Hierarchical ORAM for Interactive Two-Party Computation},
 pdfauthor={Technical design and cost study}}
\setlength{\parindent}{0pt}
\setlength{\parskip}{5pt}
\setlength{\emergencystretch}{2em}
\setlength{\headheight}{14pt}
\setlist{nosep,leftmargin=1.6em}
\pagestyle{fancy}\fancyhf{}
\fancyhead[L]{\small\textcolor{navy}{Hierarchical ORAM for 2PC}}
\fancyhead[R]{\small Construction and cost analysis}
\fancyfoot[C]{\thepage}
\newcommand{\xor}{\mathbin{\oplus}}
\newcommand{\AND}{\mathbin{\wedge}}
\newcommand{\ORC}{\operatorname{ORCompact}}
\newcommand{\PRF}{\operatorname{PRF}}
\newcommand{\occ}{\operatorname{occ}}
\newcommand{\pres}{\operatorname{present}}
\newcommand{\ceil}[1]{\left\lceil #1\right\rceil}
\newcommand{\floor}[1]{\left\lfloor #1\right\rfloor}
\newcommand{\todo}[1]{\par\textbf{TODO.} #1\par}
\begin{document}
{\color{teal}\large TECHNICAL REPORT\par}
{\color{navy}\LARGE\bfseries Binary Hierarchical ORAM\\for Interactive Two-Party Computation\par}
\vspace{5pt}
\textbf{Large routing buckets, small lookup buckets, and packed payloads}\par
17 October 2026 \quad $\cdot$ \quad Semi-honest two-server setting
\vspace{6pt}

\begin{abstract}
We adapt the hierarchical design of Microrama to streamed, interactive
two-party computation. A read invalidates its old copy and reinserts one
batch representative, eliminating key deduplication during rebuilding.
The hierarchy is binary, its base is a hash table, and no recursive position
map or deamortized schedule is required. Secret PRFs compress application
keys to 128 bits and generate all table addresses.

Rebuilding first routes records through 512-slot buckets, then refines each
bucket into 64- or 128-slot lookup buckets using recursive ORCompact.
A metadata-only planner identifies exact final-bucket excess before entries
move. One 64-slot residual buffer per parent collects this excess; a further
ORCompact merges all parent buffers into a fixed table stash. Every read
searches its selected small bucket and that stash. We include the cost of
re-admitting the stash at the next rebuild.

With a batch of 1024, 64-slot lookup buckets use approximately 21.4 and
28.5 MB per lookup at $N=2^{28}$ and $2^{32}$, including rebuilding, reads,
initial PRF evaluation, duplicate handling, and packed-payload movement.
Rebuilding accounts for about 11.1 and 15.0 MB and is now the larger component.
The current Circuit OPRAM planning ledger uses 226.5 and 286.9 MB.
These are communication estimates for specified kernels, not measured
throughput. Payloads use a concrete XOR-share/word-OT backend; base-OT setup,
CPU, initialization, replication, and implementation validation remain separate.
\end{abstract}
%%SUMMARY%%
All traffic tables use decimal MB and omit packet framing unless stated.
The probability analysis assumes ideal fresh PRFs and treats 128-bit
fingerprint collisions separately. Its numerical evaluations are not
interval-certified proofs.

\clearpage
{\small\setlength{\parskip}{0pt}\tableofcontents}
\clearpage

\section{Scope, result, and assumptions}
This report specifies a private dictionary for fixed-length 32-byte keys and
32-byte values. Garbler $G$ and evaluator $E$ are noncolluding servers; a
static semi-honest adversary corrupts at most one of them. A query arrives as
two XOR shares of its application key. The answer and its presence bit remain
shared. Database contents are similarly shared or encoded during initialization.
The main workload is a fixed dictionary with application reads; each read
still mutates the ORAM representation. The reference service rate is
10,000 lookups/s, but no particular deployment is assumed to achieve it.

Our principal conclusion is a communication feasibility result. Refining
large routing buckets removes much of the earlier read-scan overhead while
adding modest amortized work to rebuilding. A 64-slot leaf performs best
in the evaluated average-traffic cases. A 128-slot leaf has a much smaller
largest-table stash, which can matter for latency and large-table contention.
The comparison is against the current Circuit-OPRAM-inspired compilation
ledger, not a lower bound for all Circuit OPRAM implementations.

\begin{center}\small
\begin{tabularx}{\linewidth}{lX}\toprule
Symbol & Meaning\\\midrule
$N$ & Dictionary capacity, $2^{28}$ or $2^{32}$ genuine records\\
$b$ & Request batch size; primary choice 1024\\
$n_i=b2^i$ & Logical capacity of lower level $i$\\
$L=\log_2(N/b)$ & Number of lower levels; permanent top has capacity $N$\\
$Z=512$ & Routing-bucket capacity; nominal mean occupancy 256\\
$z\in\{64,128\}$ & Lookup-bucket capacity; nominal mean occupancy $z/2$\\
$d=\log_2(Z/z)$ & Additional PRF bits for the local refinement\\
$s=64$ & Temporary residual capacity per 512-slot parent\\
$G_z(n)$ & Fixed capacity of the merged residual stash of an $n$-record table\\
$Q=2^{32}$ & Access horizon used for capacity calculations\\
$c_V=160$ & Bytes per secret swap or product of a complete 256-bit payload\\
\bottomrule\end{tabularx}
\end{center}

Public information includes capacities, batch size, level occupancy,
rebuild times, selected physical bucket indices, and message sizes. Keys,
record occupancy, absence markers, hit locations, routing controls, and
residual counts remain private. All operations have fixed public capacity.
Failure events are explicitly bounded; an implementation must handle an
overflow without silently discarding a valid entry.

\section{Hierarchy, records, and the binary rebuild schedule}
\subsection{Mutable tables and unique occupied copies}
Lower levels have capacities $b,2b,\ldots,N/2$, with a permanent top of
capacity $N$. A table of logical capacity $n$ has $2n$ main slots. After
refinement these slots form $2n/z$ lookup buckets. A separate $G_z(n)$-slot
stash holds excess records. Unlike Microrama's read-only tables, access
clears an old occupied bit before reinsertion~\cite{microrama}.

An entry consists of a 128-bit canonical fingerprint, two Boolean flags,
and a packed 256-bit value. The flags are distinct:
\begin{itemize}
\item $\occ=0$ denotes a dummy or an invalidated old copy.
\item $\occ=1,\pres=1$ denotes a present dictionary record.
\item $\occ=1,\pres=0$ denotes a keyed cached absence, also called a tombstone.
\end{itemize}
The steady metadata width is therefore $R=130$ bits. Routing tags exist
only during rebuilding. The construction maintains one occupied copy of
each represented fingerprint across main buckets and stashes.

On a hit, including a cached absence, clear the old occupied bit and carry
the entry into the new batch. On a miss, create a keyed cached absence.
Repeated misses must find that absence before returning to older table
hashes. Inserting an unkeyed dummy on a miss would reveal repeated queries.
Only the full top refresh can discard cached absences: it retires all old
hashes and empties every lower level. Application deletion and concurrent
writes require an additional functionality specification.

\subsection{One hash-table base and one batch deduplication}
The base is a hash table of logical capacity $b$, not a dense online stash.
For $b=1024$ it initially has four 512-slot routing buckets and ends with
32 buckets of 64 slots or 16 buckets of 128 slots, plus its residual stash.

At batch entry, compute canonical fingerprints, bitonically sort them, and
privately mark one representative per distinct key. Duplicate request slots
remain in the fixed batch and issue dummy probes. Representatives alone
are reinserted as occupied entries. At batch exit, a segmented propagation
copies the representative's answer through its equal-key run, then saved
sort controls restore original request order. None of these controls is
opened. This is a one-time batch cost; there is no position-map recursion.

\subsection{Direct binary-counter rebuilds}
After a batch, rebuild directly into the smallest empty lower level, merging
the new batch and every occupied smaller level. If all lower levels are full,
refresh the top together with all lower levels. Intermediate carry tables
need not be materialized. The schedule depends on batch count even when a
batch has fewer than $b$ distinct keys.

Over a full counter cycle, lower capacity $n_i$ is the destination once per
$2n_i$ accesses; the top is rebuilt once per $N$ accesses. Consequently,
\begin{equation}\label{eq:amort}
 \overline C=\sum_{i=0}^{L-1}\frac{C(n_i)}{2n_i}
                  +\frac{C_{\rm top}(N)}N.
\end{equation}
Each lower table is present half the time. The mean number of searched
tables is $L/2+1$, while the maximum is $L+1$: 10/19 at $N=2^{28}$ and
12/23 at $N=2^{32}$ for $b=1024$.

Independent ORAM replicas can service requests while one replica rebuilds.
This avoids requiring a deamortized construction, but does not remove
amortized rebuilding work or the storage cost of replicas. Static application
data allow replicas to evolve independently; application writes require a
separate consistency protocol.

\section{PRFs and the mixed representation}
\subsection{All addresses come from secret PRFs}
Compute the 16-byte canonical key using AES-CMAC under a jointly secret key.
For a fixed 32-byte application key, this uses two AES encryptions. Use fresh
independent secret contexts for each published table epoch's real and dummy
addresses. A table's AES output supplies both its parent-bucket bits and
the $d$ refinement bits. This is one PRF evaluation per retained rebuild
slot, not one evaluation per routing stage. Dummy-slot evaluations are
also charged to preserve the fixed circuit schedule.

We use the NIST-listed 5120-AND AES encryption datapath and 1280-AND key
schedule~\cite{nistcircuits}. Thus initial key compression costs 10,240
ANDs per query. The shared CMAC subkey and expanded-key setup cost 6400
ANDs once per initialized state. A fresh table's two key schedules are
charged at 2560 ANDs, with two 128-bit shared key inputs. All Boolean
ANDs cost 32 bytes in the half-gate ledger~\cite{halfgates}. XORs contribute
CPU work but no garbled-table bytes.

For a read, the ledger computes both a true address and a fresh dummy
address, then privately selects between them. The dummy PRF input contains
a unique epoch/request identifier. Neither party knows the unused dummy
address in clear, preventing identification of the hit level by comparison.
The next table's selected index is disclosed only after updating the hidden
found flag. Requested bucket indices are visible; fingerprints and full
routing tags remain secret.

\subsection{Boolean metadata and packed value carriers}
Metadata stays in ordinary Boolean GC representation. Each party retains
16 bytes per encoded metadata bit. A 256-bit payload is represented as
ordinary XOR shares $V=V_G\xor V_E$, each share occupying two 128-bit
words. This is a precise packed-value interface; it is not an assertion
that two ordinary Boolean wire labels encode a 256-bit value.

Payloads stay in this representation throughout compaction, routing,
selection, and result propagation. No share-to-GC re-input is needed for
pure movement. A later application that computes on the returned value
must pay for its selected representation boundary. Initialization must
establish these payload shares; conversion from a different preexisting
opaque encoding is not free.

\subsection{A 160-byte secret word operation}
Let the GC output control bit $c$ have zero-label color $c_G$ and active-label
color $c_E$. With the standard Free-XOR offset satisfying
$\operatorname{color}(\Delta)=1$, these give local XOR shares
$c=c_G\xor c_E$, with no network
conversion. For payload shares $x_P,y_P$, define $d_P=x_P\xor y_P$.
Use the two word OTs below with fresh independent random 256-bit masks:
\begin{align*}
 G\longrightarrow E &: (r_G,\ r_G\xor d_G),\quad\text{choice }c_E,
                   \quad\text{received }u_E;\\
 E\longrightarrow G &: (r_E,\ r_E\xor d_E),\quad\text{choice }c_G,
                   \quad\text{received }u_G.
\end{align*}
Set
\begin{align*}
 t_G&=c_Gd_G\xor r_G\xor u_G,&
 t_E&=c_Ed_E\xor r_E\xor u_E,\\
 x'_P&=x_P\xor t_P,& y'_P&=y_P\xor t_P.
\end{align*}
Then $t_G\xor t_E=c(x\xor y)$, so the outputs are a conditional swap.
Replacing $d$ by a single word gives a secret bit-times-word product,
used for one-hot read extraction.

With chosen-message OT extension, each directed OT budgets two 32-byte
ciphertexts and 16 bytes of extension matrix communication at security
parameter 128: 80 bytes per direction, 160 bytes total~\cite{iknp}.
Base OTs are amortized separately. The two directions have independent
extension instances and domain-separated pad hashes. Fresh masks hide each
output share. The circuit controls remain shared, and no routing schedule
is opened. This explicit semi-honest construction requires implementation
validation; malicious security would require additional checks.

\section{A complete rebuild with local refinement}\label{sec:rebuild}
\subsection{Source compaction includes old stashes}
A lower rebuild into capacity $n$ initially has $2n$ \emph{main-array}
source slots: all smaller main tables plus the new batch padded to $2b$.
A top rebuild has $4N$ main-array source slots. Old global stashes add
physical slots beyond these counts and cannot be ignored.

We use two compactions. First compact the main arrays by occupied bits and
keep the first $n$ slots. At the top, first clear $\occ$ for cached absences.
There are at most $n$ occupied retained keys in all sources combined, so
keeping this main-array prefix is safe. Append every source stash at fixed
public offsets after slot $n$, clear absent records there at the top, and
pad the combined array to $2n$. A second ORCompact keeps its first $n$
slots. The sum of source stash capacities is at most $n$ in every evaluated
configuration, and the model checks this assertion. The base has no old
stash and needs only the first compaction.

This extra pass avoids a hidden source-size increase or a free stash merge.
All old tags can be discarded before these compactions. Next compute the
new PRF tags for the $n$ retained slots and distribute these slots publicly
into $2n/Z$ parents with $Z/2$ designated input positions each, padding
the rest. Every initial parent has at most $Z/2$ occupied slots.

\subsection{The 512-slot butterfly}
Let $p=2n/Z=2^h$. A binary butterfly routes records by $h$ parent-index
bits. It has $h$ stages, each containing $p/2$ independent merge-splits.
Every merge-split acts on $2Z=1024$ slots. It marks real entries for the
left output and enough dummies to make exactly $Z$ marks, then invokes
ORCompact~\cite{orcompact}. This works unless a real side exceeds $Z$;
Section~\ref{sec:failure} bounds that event.

The $d$ extra leaf-index bits travel through this network too. They are
not used to choose a 512-slot parent but are needed for local refinement.
Dropping them and silently recomputing them would omit another PRF cost.
After global routing, the parent-index bits can be discarded, leaving
130 metadata bits plus $d$ suffix bits per entry.

\subsection{Select exact leaf excess using only metadata}\label{sec:selection}
A naive shrinking binary split can eject records at an intermediate node
even when their final leaf has room. We instead determine final leaf excess
first, using one-bit planners, and then recursively split a feasible array.
For one 512-slot parent, let $g=512/z=2^d$ and let $X_j$ count occupied
records whose suffix names leaf $j$.

For every leaf $j$, perform the following fixed procedure:
\begin{enumerate}
\item Form $M_j[i]=\occ_i\AND[\mathrm{suffix}_i=j]$ for all 512 slots.
\item Compute the ORCompact routing controls for $M_j$.
\item Run those controls \emph{backwards} on the public one-bit vector
      $1^z0^{512-z}$. This maps the first $z$ output positions back to
      selected input positions.
\item AND the inverse mask with $M_j$. This keeps exactly
      $\min(X_j,z)$ marked entries of leaf $j$.
\end{enumerate}
If $X_j<z$, the inverse mask also contains unmarked slots; the last AND
removes them. ORCompact is stable on marked entries, although any fixed
private choice of $z$ entries would suffice. Invert the keep mask within
$M_j$ to obtain that leaf's excess mask. These masks are disjoint across
leaves, so their union is a free XOR. The final residual count is exactly
\begin{equation}\label{eq:residual}
 r=\sum_{j=0}^{g-1}(X_j-z)_+.
\end{equation}
Only one-bit arrays move during this selection step. Neither keys nor
payloads are routed separately for each leaf.

\subsection{One local residual stash and recursive ORCompact}
ORCompact the complete 512-slot parent once using the residual mask. The
first $r$ records are residuals. Copy the first 64 slots' data into a local
stash at fixed public offsets, set its occupied bits to $[i<r]$, and set
the main occupied bits to $\occ_i\AND\neg[i<r]$ for $0\le i<64$.
Slots $r,\ldots,63$ keep their original occupied bits; the remaining main
slots are also unchanged. The copies are unconditional and
local; no conditional payload move is added. Preserve the presence bit in
the stash. Check $r>64$ privately as an overflow condition.

Every remaining leaf now has at most $z$ occupied records. Recursively
partition the main 512 slots into two 256-slot arrays, then four 128-slot
arrays, and, for $z=64$, eight 64-slot arrays. Each node uses the same
exact-padding ORCompact split. No intermediate node can overflow: its
occupied count is at most the sum of the capacities of its descendant
leaves. Consumed suffix bits may be dropped after each stage.

This construction uses \emph{one} shared local stash per original parent,
not a stash per split node. Its capacity depends on final excess only.
In the worst deterministic input, 512 keys can name one leaf and require
$512-z$ residual slots. The choice of 64 is probabilistic and is justified
over the access horizon below.

\subsection{Merge all parent residuals}
There are $p=2n/512$ parents. Concatenate their fixed 64-slot buffers and
ORCompact the resulting
\begin{equation}
 p\cdot64=n/4
\end{equation}
slots by occupied bits. Keep the first $G_z(n)$ slots as the table's
persistent global stash, preserving all 130 metadata bits and packed
payload shares. Check whether more occupied entries remain. Because $n/4$
is already a power of two, this collector needs no extra padding.
The local buffers are temporary and can be reclaimed after publication.

The table now consists of $2n/z$ small lookup buckets and one global stash.
The total occupied count remains at most $n$. Entry uniqueness holds across
both parts. No residual count, destination, or compaction control is revealed.

\section{Lookup, deletion, and batch parallelism}
For each occupied table, first select the request's true or dummy index
using the hidden found flag. Reveal only that selected index. The evaluator
requests a fresh circuit for the corresponding $z$-slot bucket and the
table's fixed $G_z(n)$-slot stash. Both parties already know the table
version and the public set of requested indices.

Inside the circuit, compare the 128-bit fingerprint with every occupied
candidate, gated by $\mathrm{active}=\mathrm{representative}\AND\neg\mathrm{found}$.
A dummy physical probe therefore cannot delete an occupied record even
if its stored key happens to match. A present
record returns its value; a cached absence stops further real probes and
returns absent. Compute all bit-times-payload products in parallel and
XOR their output shares. No sequential chain of payload OTs is required
for scanning a bucket or stash. Clear the occupied bit of each matched
entry. Unique occupied copies and request deduplication make the matching
requests disjoint.

All requests to the same public bucket must use one combined circuit with
aggregated deletion masks, or serial fresh circuits linked through the
updated state. The global stash is shared by \emph{all} requests to that
table, so its updates must likewise be combined. Independent circuits
reading the same old stash version could restore records deleted by another
request. The model retains $bG_z(n)$ stash comparisons per batch and does
not assume this coordination removes them.

Garbled tables, gate tweaks, and mutable output labels must be fresh for
each invocation/version. On-demand garbling solves the scheduling problem
of knowing which bucket circuits are needed. It does not authorize
reusing a one-time garbled circuit on changing bucket contents.

Queries and buckets within one table stage parallelize, as do the
merge-splits within a rebuild stage. A single query's table probes retain
a dependency: the next disclosed index depends on whether it has already
found its key. Computing candidate PRFs ahead of time does not allow all
real indices to be disclosed at once. This limitation is also present in
Microrama's lookup description~\cite{microrama}.

\section{Overflow and fingerprint error budgets}\label{sec:failure}
\subsection{512-slot routing buckets}
After balanced initial placement, an intermediate bucket receives a sum
of fresh-PRF routing indicators with mean at most $Z/2$. Capping previous
buckets only removes records, so a first-overflow argument compares with
the uncapped network. Chernoff gives
\begin{equation}
 \Pr[X\ge Z]\le(e/4)^{Z/2}.
\end{equation}
At $Z=512$ this is at most $2^{-142.67}$ per bucket-stage. Multiply by
the public number of bucket-stage tests over the horizon. For the primary
cases the resulting bounds are below $2^{-111}$.

\subsection{A 64-slot local residual buffer}
For a parent with $g=512/z$ children, let $R=\sum_j(X_j-z)_+$. If $R>s$,
some subset of $a$ overflowing leaves jointly contains at least
$u_a=az+s+1$ records. Before conditioning on the parent load, its mean
is at most $\mu_a=az/2$. A subset union bound gives
\begin{equation}\label{eq:localtail}
 \Pr[R>s,\ X_{\rm parent}\le512]
 \le \sum_{\substack{1\le a\le g\\u_a\le512}}
 \binom ga\exp\{-u_a\ln(u_a/\mu_a)+u_a-\mu_a\}.
\end{equation}
Terms with $u_a>512$ are impossible when the parent routing succeeds.
This is not a claim conditioned on an adversarially fixed full parent;
it is a joint event bound under the fresh hash distribution. Parent-routing
failure is budgeted separately.

The number of parent refinements over $Q$ accesses is
\begin{equation}
 H=\frac{Q(L+2)}{512}.
\end{equation}
Using $s=64$, Equation~\eqref{eq:localtail} and a union bound yield lifetime
probabilities below $2^{-87}$ for 64-slot leaves and below $2^{-91}$ for
128-slot leaves in both primary cases.

\subsection{The merged table stash}
For a table of capacity $n$, let $M=2n/z$ and $Y\sim\operatorname{Pois}(z/2)$.
The exact moment generating function of a leaf's excess is
\begin{align}
 f_z(t)&=\mathbb E[e^{t(Y-z)_+}]\nonumber\\
 &=\Pr[Y\le z]+e^{-tz+(z/2)(e^t-1)}
        \Pr[\operatorname{Pois}((z/2)e^t)>z].\label{eq:mgf}
\end{align}
Under Poissonization the bucket occupancies are independent. Overflow is
monotone in the total number of balls, and
$\Pr[\operatorname{Pois}(n)\ge n]\ge1/2$ for integer $n$. Coupling with
exactly $n$ balls therefore gives
\begin{equation}\label{eq:globaltail}
 \Pr[R_{\rm table}>G]
 \le2\inf_{t>0}\exp[-t(G+1)]\,f_z(t)^M.
\end{equation}
Fewer occupied records only decrease the event. This argument does not
assume independence between levels or rebuilds.

There are $Q/b=2^{22}$ rebuilds at $b=1024$. Choose each $G_z(n)$ so that
Equation~\eqref{eq:globaltail} is below $2^{-88}$, then round up to a
multiple of eight. A union bound gives a global-stash lifetime budget
below $2^{-66}$. The actual evaluated sums are smaller. The optimizer
only chooses a positive Chernoff witness $t$; evaluating any such witness
already gives a valid analytic bound.
%%CAPACITIES%%

At the top, the Poisson model predicts about 3.20/51.22 residual records
for 64-slot leaves at $N=2^{28}/2^{32}$. For 128-slot leaves the corresponding
expectations are approximately $5.1\cdot10^{-6}/8.2\cdot10^{-5}$. Expected
size does not determine a safe capacity: rare large excesses still require
a stash. In particular, omitting residual handling for 128-slot leaves
does not meet the stated lifetime target.

\subsection{Finite fingerprints and implementation qualification}
The above are routing-capacity bounds under ideal PRFs and distinct
fingerprints. With $D$ distinct application keys, 128-bit compression has
the elementary collision bound $\binom D2/2^{128}$. For $N=2^{32}$ and
another $2^{32}$ distinct queried keys, $D\le2^{33}$ gives a bound near
$2^{-63}$. The routing bounds therefore do not establish a combined
$2^{-64}$ failure guarantee for this maximum workload. A 128-bit label
length also does not imply a 128-bit service-lifetime error guarantee.

Fresh table keys are chosen after the source record set is fixed. This
supports a conditional analysis for adaptive access histories; a complete
reduction still needs to specify exposure of truncated outputs, the PRF
instantiation, initialization, and the canonical-key collision event.
\todo{If a strict total lifetime target is required, allocate its budget
between fingerprint collisions, finite-PRF advantages, and capacity failures;
consider a wider fingerprint or an explicit collision-resolution mechanism.}
\todo{Certify the numerical Chernoff witnesses with interval arithmetic.
The present calculation uses stable floating-point tail and log-MGF formulas
and rounds capacities upward, but is not machine-checked analysis.}

\section{Circuit and communication accounting}\label{sec:cost}
\subsection{ORCompact and exact-padding planners}
For $a=2^k$ slots, ORCompact has exactly
\begin{equation}
 S(a)=\frac a2\log_2a
\end{equation}
conditional swaps~\cite{orcompact}. A constructive Boolean schedule planner
used in this ledger costs at most
\begin{equation}
 P(a)=a(k-\tfrac12)-k+1
\end{equation}
ANDs. To split $a$ slots evenly, forming exactly $a/2$ marks adds at most
$3ka+4a-2k-2$ ANDs, giving
\begin{equation}
 A_{\rm split}(a)=4ka+\tfrac72a-3k-1.
\end{equation}
These Boolean counts are this report's compilation allowances, not gate
counts claimed by the ORCompact paper. Appendix~\ref{app:planner} explains
the decomposition and the dummy-rank width requirement.

A conditional swap of an $r$-bit Boolean record costs $r$ ANDs: compute
$t=c\AND(x\xor y)$ bitwise, then $x'=x\xor t,y'=y\xor t$.
The complete mixed-record operation therefore costs $32r+c_V$ bytes.
Public copies, public permutations, and XOR-share fan-out add no garbled
tables, although they do incur memory and CPU work.

For the global butterfly, $p=2n/512$ and $h=\log_2p$. Its $ph/2$ calls
each have $S(1024)=5120$ record swaps. The model conservatively carries
all $h+d$ routing bits through this global network, for
\begin{equation}
 C_{\rm butterfly}=
 \frac{nh}{512}\left[32A_{\rm split}(1024)
              +5120\{32(130+h+d)+c_V\}\right].
\end{equation}
The initial PRF cost is $5120n$ ANDs, plus table-key setup. Old source
compactions and the global residual collector move only 130 metadata bits.

\subsection{Local refinement costs}
Each leaf's excess planner costs
$512d+P(512)+S(512)+512$ ANDs: membership flags, schedule planning,
one-bit inverse routing, and mask intersection. Multiply by $512/z$.
The residual extraction uses $S(512)=2304$ complete-entry swaps and
$P(512)+197$ planner/check ANDs. Recursive splits cost the sum of the
node costs at sizes 512, 256, and optionally 128; already consumed suffix
bits are discarded. The resulting kernel is:
%%LOCALCOUNTS%%

For each rebuilt table, multiply these counts by $p=2n/512$. The global
collector adds $P(n/4)$ planner ANDs and $S(n/4)$ swaps of a 130-bit
metadata record plus its packed value. The two-stage source admission
adds a full second $2n$-slot compaction whenever old stashes exist.
Both costs are included before applying Equation~\eqref{eq:amort}.

\subsection{Read and frontend ledger}
One occupied table scans $z+G_z(n)$ slots. A conservative read allowance is
\begin{equation}\label{eq:read}
 A_{\rm read}(n)=10240+134[z+G_z(n)]+h+d+4
\end{equation}
ANDs, plus $z+G_z(n)$ packed word products at $c_V$ bytes each. The
10240 term charges true and dummy PRFs. The scan coefficient covers
fingerprint equality, occupancy/active gating, presence reduction, deletion,
and accumulation. Disclosed address bits add a small explicit bidirectional
term. Sum lower levels with weight $1/2$ and the top with weight one.

For $b=2^\beta$, a bitonic network has
$b\beta(\beta+1)/4$ comparators. The model charges 511 ANDs per comparator
for 128-bit comparison, fingerprint swapping, and inverse presence-bit
routing, plus one packed result swap on the inverse network. It also
reserves 140 ANDs and three packed selectors per adjacent input for a
work-efficient parallel segmented propagation. This is deliberately more
than the earlier sequential propagation allowance.

The initial 32-byte secret-shared query input costs 16,384 bytes under
the existing input convention; AES-CMAC adds 327,680 bytes. No initial
database hashing, base-OT setup, client delivery, or application computation
on the result is amortized into the steady-state rows.

\section{Numerical results and sensitivities}\label{sec:results}
\subsection{Primary traffic comparison}
%%RESULTS%%
All rows use complete application lookups. Our hierarchy directly stores
canonical keys and has no recursive position maps or outer pair of cuckoo
table probes. The Circuit OPRAM ledger includes both cuckoo probes and its
recursive machinery. It also uses ordinary Boolean payload movement in
its primary row. These are complete alternative layouts, not identical
circuits with only one kernel substituted. Its selected batch is 32;
the hierarchical primary batch is 1024. The comparison does not equalize
batch-fill delay or measured latency.

Refinement makes rebuilding the larger component for both leaf sizes.
The 64-slot choice minimizes average traffic in the primary rows, but at
$N=2^{32}$ its top lookup scans $64+264=328$ slots, versus $128+64=192$
for 128-slot leaves. Thus the smaller main bucket is not universally the
smaller read operation. Much of the average advantage comes from the
smaller tables' more modest stashes.

\subsection{Rebuild decomposition}
%%DECOMPOSITION%%
The dominant rebuild terms move Boolean metadata through the global
butterfly and source compactions. One-bit planners are important but do not
make record movement free. Packing the payload does not pack the 128-bit
fingerprint or the routing bits used by GC.

\subsection{Batch-size sensitivity}
%%BATCHSWEEP%%
Larger batches reduce the number of levels and the per-lookup rebuild
frequency of the smallest tables, while increasing initial sort work and
batch-fill time. At rate $q$, ideal mean fill delay is $(b-1)/(2q)$:
51.15 ms at $b=1024$ and 10,000 requests/s. This is only fill delay,
before circuit service and queueing. The evaluated sweep is not a full
latency-constrained optimization.

\subsection{Packed operation sensitivity}
%%PAYLOADSWEEP%%
The 160-byte row corresponds to the specified two-direction word-OT
backend. The 80/320-byte rows are parameter sensitivities. The 8192-byte
row is an ordinary 256-AND Boolean-payload movement reference, with the
same metadata kernels; it is not the storage cost of two packed shares.
It illustrates how much of the conclusion survives without a packed
movement discount.

\section{Space, parallel work, and latency}\label{sec:resources}
\subsection{Persistent state and rebuilding workspace}
Before a full top rebuild, the maximum main-array allocation is
$4N-2b$ slots. Add the public global stash capacity of every table. A
packed record uses $16\cdot130+32=2112$ bytes per party. This gives:
%%RESOURCES%%
These are steady representation allocations, not peak process memory.
Rebuilding also needs transient tag labels, saved secret controls, metadata
planner workspace, and the $n/4$-slot parent residual collection. A copied
main-array source or a second compaction buffer must be charged if the
implementation does not operate in place. OT seeds, queues, GC streaming
buffers, and extra replicas are additional. The model does not yet price
storage I/O or claim that a whole state fits in one machine.

\subsection{Mixed-protocol depth}
GC metadata planning can be computed before payload routing, because the
planning depends only on keys, flags, and tags. The packed word operations
remain interactive: an OT layer's sender inputs depend on the payload
shares produced by previous layers. Independent swaps in a layer batch
together, but precomputing random OTs does not remove all online dependencies.

The primary top rebuild has 512-bucket butterfly depth $10h$ payload
swap layers, plus source compactions, local extraction/refinement, and
global stash collection. The model reports this serial-layer count; it is
not a measured number of RTTs because OT preprocessing and transport
scheduling can differ. Read payload products for all candidate slots run
in parallel, so scanning a 64- or 128-slot leaf does not imply that many
serial OT rounds. Duplicate propagation uses a parallel prefix rather
than a linear chain in the new ledger.

\subsection{Replicas and throughput}
At 10,000 lookups/s, multiply the unframed MB/query by 10 to get aggregate
GB/s across both communication directions. Directional provisioning must
separate predominantly garbler-to-evaluator GC tables from the symmetric
word-OT traffic. Aggregate bandwidth alone does not establish a state's
batch service time or the replica count needed while the top rebuilds.
For batch service time $T_b$, a steady-state group count must satisfy at
least the relevant arrival/service constraint; long periodic rebuilds need
an explicit availability schedule as well. Neither annual cloud cost nor
CPU-limited throughput is claimed here.

\subsection{An optional parent-local persistent stash}
Instead of globally merging the residuals, one could retain each parent's
64-slot stash and search it with that parent's small bucket. The parent
index is already determined by the disclosed small-bucket address, so this
does not add address leakage if both addresses derive from the same
true-or-dummy choice. It scans $z+64$ slots and avoids the global collector,
but adds $n/4$ persistent slots per table. Different children share a parent
stash and must coordinate updates. Source admission must include those
extra slots. This alternative is not included in the primary numbers,
which follow the requested merged-residual organization.

\section{Validation, remaining work, and recommendation}
The cost program is a ledger of explicit kernels rather than a full 2PC
implementation. Its checks cover binary-counter frequencies, exact
ORCompact swap arithmetic, final-excess selection, conservation of keys,
recursive leaf capacities, threshold boundary cases, and numerical tail
formulas. Constructive simulations include empty, balanced, highly skewed,
and random parents, plus residual counts 0, 1, 63, 64, and 65. Overflow
conditions are tested before truncation.

\textbf{Recommendation.} Pursue the 512-to-64 design as the primary
communication candidate, retaining 512-to-128 as a useful alternative
when largest-table stash scanning or contention matters. The additional
refinement and stash work is comfortably below the savings in ordinary
512-slot reads. Rebuilding again becomes the main amortized traffic term,
as anticipated. No early rejection against the current Circuit OPRAM
ledger is justified by these counts.

The following work remains before an implementation or stronger performance
claim:
\begin{enumerate}
\item \textbf{Compile and count the actual circuits.} Replace the planner,
      scan, and parallel-propagation allowances by emitted-circuit counts;
      prove their label interfaces compose safely.
\item \textbf{Implement the packed backend.} Validate GC color-bit sharing,
      both OT directions, fresh masks/tweaks, scheduling, and application
      input/output boundaries. Benchmark local hashing and XOR work.
\item \textbf{Complete the security argument.} Prove adaptive lookup
      indistinguishability with absent tombstones, grouped bucket/stash
      updates, fresh epochs, and bounded failures. The sketch here is not
      a complete simulation proof.
\item \textbf{Certify probability calculations.} Check the analytic proof
      conditions and numerical witnesses, and choose a total fingerprint,
      PRF, and capacity failure budget for the intended service lifetime.
\item \textbf{Measure a distributed rebuild.} Account for metadata-label
      traffic inside each party, peak memory, SSD/RAM I/O, OT dependency
      layers, and the largest rebuild pause.
\item \textbf{Choose a deployment schedule.} Choose batch size and replica
      count from measured service time and latency targets; application
      writes additionally require replica consistency.
\end{enumerate}

\appendix
\section{Planner derivation and finite counts}\label{app:planner}
For $a=2^k$, a concrete ORCompact planner comprises balanced subtree
popcounts, offset additions, shared threshold decoding, and leaf controls.
Their AND allowances are respectively
\[
 2a-2k-2,\qquad a-k-1,\qquad a(k-4)+2k+4,\qquad a/2.
\]
Their sum is $P(a)=a(k-1/2)-k+1$. A shared decoder for all comparisons
of a $q$-bit secret threshold against public indices uses
$D(q)=2^{q+1}-2q-2$ ANDs. This avoids charging a separate full comparator
for every index.

Exact-padding marks first identify real-left entries and their total,
then compute the needed dummy quota. Use \emph{exclusive} dummy prefix
ranks, which are at most $a-1$, and select a dummy iff its rank is below
the quota. Inclusive ranks can equal $a$ and would require another bit.
A work-efficient parallel prefix scan, comparisons, and gating give the
$3ka+4a-2k-2$ marking allowance. The construction assumes both real
output counts fit; the probability argument or the preselection invariant
establishes that condition.

For extracting a parent's residuals, $P(512)=4344$ and $S(512)=2304$.
The additional 197 planner/check ANDs cover a nine-bit residual count,
the 64 prefix occupancy flags, and the $r>64$ test. At most $512-z$
records can be residual, so nine bits suffice. The data copy into fixed
stash positions is public; live flags preserve its secrecy.

The complete local kernel counts and per-level rebuild contributions are
stored in the accompanying JSON. The model retains all global routing
bits through the butterfly as a conservative choice; discarding consumed
bits there offers a further optimization that is not included.

\section{Reproduction and source provenance}
From the repository root, with Python, SciPy, and mpmath available:
\begin{verbatim}
python -m mpc_cost.hierarchical_split_cost
python -m unittest mpc_cost.test_hierarchical_split -v
python -m mpc_cost.render_hierarchical_split_report
\end{verbatim}
Compile the delivered self-contained LaTeX source twice with pdfLaTeX.
It embeds all numerical tables and bibliography entries and requires no
project-specific input files or external figure files. The model writes
\path{mpc_cost/results/hierarchical_split_comparison.json}; the report
generator writes \path{output/latex/hierarchical_batch_oram_report.tex}.

Microrama supplies the hierarchical/butterfly design and the necessity of
cached absences. ORCompact supplies the compaction network and exact swap
counts. The mutable 2PC adaptation, exact-excess preselection, mixed
payload interface, stash bounds, and numerical ledger are the present
analysis. The Circuit OPRAM comparison uses the repository's 16 October
2026 interactive-shuffle ledger, whose traffic is 226.506916/286.945386 MB
unframed and 237.832261/301.292656 MB with 5\% framing. That ledger is
itself a conditional circuit-composition model.

\begin{thebibliography}{9}
\bibitem{microrama} Anonymous submission.
\emph{Microrama: Microsecond Oblivious Access to Gigabytes of Data}.
User-supplied ongoing draft, \path{microrama.pdf}, especially Sections 2.1--2.4.
The draft's native TEE measurements are not used as 2PC performance measurements.
\bibitem{orcompact} Sajin Sasy, Aaron Johnson, and Ian Goldberg.
\emph{Fast Fully Oblivious Compaction and Shuffling}. CCS 2022, extended version.
\url{https://eprint.iacr.org/2022/1333}. Sections 3 and 5.1, Appendix B.
\bibitem{nistcircuits} NIST. \emph{Circuit Complexity: List of Circuits}.
AES-128 with expanded key: 5120 ANDs; with key schedule: 6400 ANDs.
\url{https://csrc.nist.gov/projects/circuit-complexity/list-of-circuits}.
Accessed 17 October 2026.
\bibitem{halfgates} Samee Zahur, Mike Rosulek, and David Evans.
\emph{Two Halves Make a Whole: Reducing Data Transfer in Garbled Circuits
Using Half Gates}. EUROCRYPT 2015.
\url{https://www.cs.virginia.edu/~evans/pubs/ec2015/}.
\bibitem{iknp} Yuval Ishai, Joe Kilian, Kobbi Nissim, and Erez Petrank.
\emph{Extending Oblivious Transfers Efficiently}. CRYPTO 2003.
\url{https://www.iacr.org/archive/crypto2003/27290145/27290145.pdf}.
The 160-byte word operation is this report's explicit chosen-message
accounting, including a conventional 128-bit extension-matrix allowance.
\bibitem{opram} T.-H. Hubert Chan and Elaine Shi.
\emph{Circuit OPRAM: Unifying Statistically and Computationally Secure
ORAMs and OPRAMs}. Cryptology ePrint Archive, 2016/1084, revised 2020.
\url{https://eprint.iacr.org/2016/1084}.
The comparison uses a repository compilation template rather than a
performance number reported in this paper.
\end{thebibliography}
\end{document}
